\documentclass[11pt,a4paper]{article}
\usepackage[T1]{fontenc}
\usepackage[utf8]{inputenc}
\usepackage{amsmath,amssymb,amsthm,newtxtext,newtxmath}
\usepackage[margin=1in]{geometry}
\usepackage{microtype,booktabs,array,graphicx,xcolor}
\usepackage[round,authoryear]{natbib}
\usepackage[hidelinks]{hyperref}
\usepackage{tikz}
\usetikzlibrary{svg.path}
\usepackage{caption}
\captionsetup{font=small,labelfont=bf}
\linespread{1.04}
\renewcommand{\topfraction}{0.9}
\renewcommand{\textfraction}{0.08}
\renewcommand{\floatpagefraction}{0.8}
\setlength{\parindent}{1.2em}
\setlength{\parskip}{2pt}
\setlength{\emergencystretch}{2em}
\newtheorem{definition}{Definition}
\newcommand{\E}{\mathbb{E}}
\newcommand{\I}{\mathbf{1}}
\hypersetup{pdftitle={Decentralized Markets for Crowdsourced Mathematics},pdfsubject={Roles, Incentives, and Institutional Design},pdfauthor={Varifold Lab}}
% MAP_MACRO_END
\begin{document}
\begin{center}
{\LARGE\bfseries Decentralized Markets for\\[3pt] Crowdsourced Mathematics}\par
\vspace{8pt}
{\large Roles, Incentives, and Institutional Design}\par
\vspace{7pt}
{\normalsize Varifold Lab}
\end{center}
\begin{abstract}
\noindent AI systems can spend computational resources to search for candidate proofs, while formal proof assistants can verify successful proofs independently of the identity of their producer. This separation between resource-intensive search and independently checkable verification creates a distinctive problem of resource allocation. Human researchers and AI agents may pursue the same mathematical objective in parallel, incurring labor and computational costs before the value of any particular line of attack is known. Traditional grants allocate resources ex ante to identified researchers and proposed projects, while prizes reward successful results ex post. Neither mechanism by itself specifies how continuously arriving information should guide funding across many distributed attempts. We develop an abstract framework for markets that connect demand for mathematical knowledge to funding, proof work, verification, and settlement. It distinguishes institutional roles, sponsorship from investment, and formal acceptance from contribution attribution. We compare support for unsuccessful exploration, rewards for intermediate results, and sharing among competing successful submissions, and examine investment under heterogeneous beliefs and coordination across platforms. Emerging initiatives illustrate how parts of this framework are already taking shape and where additional arrangements would be needed. Crypto-based settlement and prediction markets offer possible interfaces whose benefits depend on their implementation. The central question is which combinations of funding, information, and reward rules produce useful, reusable mathematical progress at a sustainable resource cost.
\end{abstract}

\section{Introduction}
Formal proof assistants provide a common means of verifying mathematical proofs. AI changes the scale and economics of producing candidates for verification. An agent can explore many proof strategies, decompose a target into lemmas, and revise its approach in response to compiler feedback, but such proof search can require substantial computation, much of which may be spent on unsuccessful attempts \citep{costquality}. Successful derivations can be checked under a specified formal environment, while the time, compute, and judgment needed to find them remain uncertain. The resulting economic problem concerns how to allocate resources before an accepted proof exists.

These developments raise a broader coordination question: how might proof search be organized when multiple human and AI contributors can work on different parts of a problem? One possible arrangement is for a sponsor to fund a mathematical target, for people and agents with different skills to pursue different approaches or subproblems, and for a proof assistant to verify the resulting derivation. Such an arrangement raises questions about how to select tasks, attract contributors, allocate computation, and reward unsuccessful searches and useful intermediate results. A grant can finance a team's exploration, and a terminal prize can express demand for a result; the contracts connecting funders, contributors, and service providers determine how support reaches individual attempts and reusable contributions.

Several recent developments address complementary parts of this coordination problem. OpenMath's tooling connects theorem tasks, local proof work, submission commitments, and reward claims \citep{openmath}. The Justin Sun Prize expresses demand through cryptocurrency-denominated mathematical awards \citep{sun}. Prove2Me organizes collaborative formalization around missions and supporting statements, while Agent Hunt studies bounty-mediated coordination among proving agents \citep{prove2me,hunt}. Other efforts address information: Kalshi's filed Millennium Prize contract concerns an institutional award event \citep{kalshi}. These examples motivate questions about the interfaces between demand, work, information, and payment.

These projects and programs are not one integrated system. Together, however, they make visible the components of an emerging market for machine-verified knowledge: funders who express demand; people and agents who supply reasoning and formalization; proof checkers that validate derivations; platforms that decompose and record work; and financial or compute infrastructure that supports participation. Mathematical proof offers a particularly clear setting for examining these relationships because accepted derivations can be checked mechanically, although the intended meaning and scholarly importance of a claim still require separate review \citep{lean}. Earlier work on blockchain-supported collaborative formalization also provides a foundation for studying how these components might connect \citep{lim}.

To analyze those connections, we define roles independently of the organizations that currently perform them. A sponsor funds an outcome it values; an investor finances work for a contingent return; a solver performs that work. Coordination, verification, information provision, and settlement each involve further decisions and obligations. One institution may combine several roles, and several institutions may compete or cooperate within a role. The organizational approach of \citet{aghiontirole} provides a useful starting point: financing, creating, and using an innovation need not reside with the same actor. This abstraction allows the analysis to accommodate different platforms and funding arrangements while keeping their responsibilities and conflicts visible.

Within this framework, we compare ways to finance exploration, reward intermediate progress, and allocate a terminal pool among competing submissions. We then examine how investors with different information assess those arrangements and how task, proof, and payment records can be reconciled across platforms. The mechanisms are design options with explicit assumptions and budget constraints. Their suitability depends on the research objective, the available evidence, and the institutions able to administer them. Existing practices provide concrete settings in which to assess those assumptions and identify missing arrangements.

Decentralization has several dimensions: participation, task selection, verification, funding, records, and governance. A proof may be checked by a public kernel while task curation, compute delivery, or prize adjudication remain centralized. Crypto-based contracts offer one possible means of committing funds and transferring claims across participants; prediction markets offer one possible source of information before a proof is found. The framework also accommodates conventional payments, grants, and shared services. The relevant question is how a particular allocation of resources and authority affects participation, research incentives, and accountability.

The relevant economics cuts across several literatures. Research on innovation incentives distinguishes encouraging experimentation from simply rewarding measured performance \citep{manso}. Contest design shows that the merits of concentrating or dividing a prize depend on costs and the designer's objective \citep{moldovanu}. Work on public access to innovation separates the incentive to produce knowledge from the terms on which others use it \citep{kremer}. Models of heterogeneous expectations and costly information caution against treating a traded price as an objectively correct forecast \citep{harrison,grossman}. These distinctions inform the comparisons below; their implications for mathematical work require models and evidence suited to this setting.

Section~\ref{sec:roles} defines the economic objects and institutional roles. Sections~\ref{sec:funding}--\ref{sec:coordination} develop the funding, investment, and coordination framework. Section~\ref{sec:implementations} relates it to institutional examples, and Section~\ref{sec:research} identifies risks and empirical research questions. The analysis supplies a common vocabulary and a set of mechanisms to examine; their effects on mathematical progress and resource allocation remain questions for comparative evaluation.

\section{Economic environment and roles}
\label{sec:roles}
We begin with the objective being funded, the evidence that counts as progress, and the decisions needed to connect work to payment. These objects provide a common basis for comparing arrangements with different institutions and degrees of decentralization.

\subsection{Objectives, tasks, and acceptable evidence}
\begin{definition}[Mission and task]
A \emph{mission} is a versioned research objective, together with a budget, an authority for interpreting that objective, and rules for recognizing progress. A \emph{task} is a contractible component of a mission: a specified deliverable, an acceptance procedure, a deadline or submission schedule, and a rule allocating payment or credit.
\end{definition}
A mission might begin with an informal claim or a desired application. Its tasks can include statement formalization, proof construction, checking, maintenance, and expert review. Changes to an agreed task statement or verification environment create a new version, whose rules must specify how existing submissions, work in progress, and prior reward commitments are treated.

For a proof task $j$, let $s_j$ denote the canonical statement and $A_j$ the admissibility policy. The policy specifies allowed axioms, dependencies, resource limits, and evidence requirements. The verification procedure $V_j$ incorporates the task's fixed formal environment; an artifact $z$ passes when $V_j(z;s_j,A_j)=1$. Acceptance means that $z$ establishes $s_j$ under the task's stated assumptions and verification rules. Verification itself consumes resources, and its cost can vary with the submitted artifact.

The task contract contains more than the predicate $V_j$. It must identify the funded reward, eligible recipients, contribution and timing rules, treatment of duplicates, disclosure obligations, and dispute authority. A valid proof may arrive after a deadline, reproduce a previously paid contribution, or belong to a team with a pre-existing sharing agreement. These are economic and attribution questions. Kernel acceptance alone cannot settle them.

\subsection{Roles and institutions}
\begin{definition}[Role]
A \emph{role} is a set of responsibilities, decision rights, and entitlements within the ecosystem. An \emph{institution} is an organization or protocol arrangement that performs one or more roles. The assignment between institutions and roles is many-to-many.
\end{definition}
The \emph{demand owner} identifies the result or capability desired and assesses its use. The \emph{sponsor} commits funds because it values that outcome, its public benefits, or the mission of supporting research. They may coincide, but need not: a philanthropic sponsor can support a community's demand. Sponsorship can take the form of grants, service procurement, milestones, or outcome prizes. Its payoff need not be a monetary return from the reward pool.

An \emph{investor} instead commits resources in exchange for a contingent financial claim. It might pay a team's compute bills for a share of that team's eventual award. It bears research and settlement risk, and evaluates the commitment according to its own information and preferences. A \emph{solver} undertakes the work. Solvers include people, agents acting under delegated authority, and mixed teams. A researcher using personal compute occupies both the solver and investor roles; a subsidized solver may carry much less financial risk. Agents require identifiable principals or account-level obligations for budgets and payouts even when their search decisions are autonomous.

The \emph{coordinator} selects or curates tasks, manages dependencies, and maintains the canonical mission state. The \emph{verifier} supplies evidence about compliance with the acceptance procedure. The \emph{record and settlement service} registers commitments, preserves provenance, reconciles accepted outcomes with reward rules, and executes or authorizes payment. These functions can be bundled for convenience, but each has a distinct source of authority. A coordinator's preference for one proof route is different from a verifier's technical judgment, and both differ from the settlement service's interpretation of payout eligibility.

Finally, the \emph{compute supplier} provides metered execution or model access. The \emph{information provider} supplies forecasts, expert assessments, cost estimates, or records of earlier attempts. An information provider need not trade; an investor need not publish its model. Governance is a cross-cutting function that assigns decision rights, specifies appeals, and constrains changes to already accepted contracts. Figure~\ref{fig:map} summarizes the principal relationships. It shows functional interfaces rather than a mandatory organizational hierarchy.

\begin{figure}[htbp]
\centering
\includegraphics[width=\linewidth]{roles_interactions.pdf}
\caption{A functional ecosystem. Each box denotes a role or closely related pair of roles, not a named institution. Arrows identify the principal flows of tasks, evidence, resources, and information. Several institutions may occupy a box, and one institution may occupy several boxes. Governance applies throughout.}
\label{fig:map}
\end{figure}

\subsection{Claims, resource units, and public knowledge}
Four objects often called tokens must be distinguished. An \emph{inference token} is a unit used to meter a model service; its price and computational implications depend on the service. A \emph{compute credit} is a contractual entitlement to specified services or a spending allowance with a provider. A \emph{payment token} is a transferable settlement asset. A \emph{reward claim} is a right to a contingent payout under a task or financing contract. A common interface may connect these objects, but it does not make them economically interchangeable.

This distinction matters for both solvency and participation. A payment balance does not reserve model capacity. An advertised reward is not necessarily escrowed. A claim on a future prize does not guarantee cash to finance work today. A sponsor can use ordinary money while a downstream platform offers restricted service credits or programmable transfers. A new native cryptocurrency is therefore an optional implementation choice. What matters for a task is that its denomination, backing, redemption conditions, and settlement obligations are intelligible to participants.

The underlying knowledge may remain openly reusable even when payment rights are private. A reward contract can compensate a contribution without granting ownership over all future uses of a theorem. \citet{kremer} studies a different institutional setting in which incentives to innovate can be combined with public access. The useful analogy here is the separation of compensation from exclusivity, not a claim that mathematical propositions should be patented. Licenses for proof code, attribution records, and private agreements over award proceeds should each state their own scope.

\section{Financing exploration and cumulative contributions}
\label{sec:funding}
\subsection{Why a terminal prize is incomplete}
A terminal prize is attractive when the final deliverable is clear and objectively checkable. It lets a sponsor specify the output without directing each participant's method. It also shifts much of the cost of unsuccessful search to participants or their financiers. That allocation can screen out weak proposals, but it can exclude capable researchers with limited liquidity and discourage work whose main benefit is to later contributors. A grant solves some financing problems while introducing others, including the difficulty of observing how funded effort is used.

The sponsor should begin with the objective it is trying to advance. Maximizing the probability of one completed proof differs from building a reusable library, improving a formalization pipeline, or learning which approaches are unpromising. Spending and submission volume are imperfect proxies for all of these. \citet{manso} shows, in a model of delegated experimentation, why tolerance of early failure can support innovation. That result motivates attention to exploration here; it does not establish that reimbursement of every failed model call is efficient. An exploration program needs a reason to regard the funded activity as informative or promising before its final outcome is known.

A practical budget can reserve separate amounts for terminal completion, intermediate results, exploratory support, and verification or administration. Let these commitments be $B^T,B^L,B^X,B^V$, with
\[
 B^T+B^L+B^X+B^V\leq B.
\]
The inequality is an accounting restriction, not a welfare theorem. Awards and subsidies must draw on specified pools, including obligations created by later task additions. A platform cannot promise unrestricted reimbursement and still rely on a fixed sponsor budget. Nor should money reserved for verified work be silently reassigned after researchers have incurred costs.

\subsection{Two ways to support unsuccessful attempts}
Under \emph{advance compute support}, a participant receives a nontransferable or otherwise restricted allowance for approved research services. The allowance can be bound to a mission, a time interval, and a set of providers. This arrangement reduces the participant's liquidity requirement and lets the sponsor set an exposure limit before work begins. Its disadvantages include dependence on supported providers, administrative selection, and incentives to exhaust a grant whose unused balance cannot be retained. Allocation can be competitive, needs-based, or partly randomized among eligible proposals; the choice affects entry and should be disclosed.

Under \emph{ex post reimbursement}, a participant initially pays and later claims a capped share of admissible expenditure. A simple rule is
\[
 S_a=\I\{a\text{ satisfies the audit policy}\}
       \min\{\overline S_a,\rho_a\widehat C_a\},
 \qquad 0\leq\rho_a\leq1,
\]
where $\widehat C_a$ is accepted expenditure for attempt $a$, $\overline S_a$ is its reserved ceiling, and $\rho_a$ is the reimbursement rate. Attempts must be authorized or otherwise covered by a preannounced allocation rule; summing unlimited individual promises would defeat the aggregate budget. This arrangement permits more choice over execution, but favors participants who can front costs and exposes them to audit and reimbursement delay.

The two arrangements can coexist. Small initial credits can support entry, with larger tranches conditional on an approved plan, reproducible artifacts, or informative updates. Reimbursement can supplement those credits for independently funded work if the contract prevents the same expenditure from being claimed twice. A subsidy is a transfer in the investor's private cost calculation; the underlying compute remains a real cost to the sponsor and to society. Comparing arrangements therefore requires both net participant costs and total resources consumed.

Neither mechanism should turn failure itself into the object being purchased. The object is an eligible attempt under a bounded research policy, which may fail to produce the target. An informative counterexample to a formalized lemma, a reproducible obstacle, or an improved search method may have value, but these need separate acceptance criteria. An agent's inability to find a proof within its budget is not evidence that the proposition is false. Failed traces should record the procedure and budget that failed, and preserve uncertainty about other approaches.

\subsection{What an effort audit can establish}
An audit faces three separate questions: was a service actually consumed, was it used on the approved task, and was the expenditure productive? Provider receipts and authenticated job records can answer part of the first question. Task identifiers, execution traces, artifact hashes, and a gateway binding approved jobs to provider requests can supply evidence for the second. None of these establishes the third. A correctly metered job can be irrelevant; a relevant job can be wasteful; a failed job can still generate useful information.

Prompt relevance is consequently a screening signal, not proof of research effort. It can be imitated by repeating the task statement around unrelated or deliberately unproductive computation. Private local runs also make complete observation difficult. A more credible subsidy system combines caps with provider-backed usage records, sampled trace review, checks on repeated submissions, and inspection of intermediate artifacts. Where a provider is itself controlled by the claimant, an invoice alone gives little assurance of an arm's-length cost. The policy should specify treatment of affiliated providers and unverified local consumption.

There is a tradeoff between audit intensity and the openness the market is intended to provide. Requiring complete prompts can disclose unpublished ideas or private data. Strict templates can suppress unconventional methods, while extensive human review can cost more than small compute grants. A reasonable starting design uses narrow budgets, minimal task-bound evidence, random audits with announced consequences, and a route to appeal. The empirical question is how much additional useful exploration these controls buy per unit of administrative and privacy cost.

\subsection{Paying for a dependency structure}
A mission can be represented by a graph of formal statements and proposed proof dependencies. A proof-sketch identifies a route from intermediate statements to the target. A checked reduction that assumes unresolved lemmas is conditional progress: the mission is complete only when all required assumptions are discharged under the approved environment. The coordinator must prevent circular certification, accidental admission of unproved statements, and changes of meaning across interfaces. A dependency graph is useful accounting information, but is not by itself evidence that the root theorem has been established.

Rewards can be attached to approved intermediate tasks before the final proof exists. Suppose the mission reserves $b_j$ for each of a set of reviewed lemmas $j$, with $\sum_j b_j\leq B^L$. A contributor who satisfies a lemma's contract receives its allocated award even if the mission subsequently stalls elsewhere. A separate terminal pool rewards completion and integration. The arrangement compensates useful partial progress and lets specialized teams participate without solving the entire mission. It also gives the sponsor a reason to scrutinize the proposed decomposition before committing money.

The principal vulnerability is endogenous task creation. If proposers are automatically paid for every lemma they invent, they can fragment easy work into many claims, restate existing results, or choose a needlessly elaborate route. Admission to the reward schedule should therefore depend on canonicalization, relevance to an approved route or broader library goal, and the remaining budget. New routes may justify new awards, but not retroactive dilution of already contracted ones. Review should also recognize that a seemingly peripheral lemma can unlock a much shorter proof than the initially preferred route.

Reuse creates an additional externality. A contributor may supply a general result that later saves many teams substantial effort. A separately funded reuse pool can reward documented adoption or curated importance, provided that its rules resist circular citation and self-generated dependencies. This is a design option rather than a requirement to pay perpetual royalties on mathematical facts. The system should distinguish scholarly credit, one-time task payment, team revenue shares, and any subsequent reuse awards. They serve different purposes and need not have the same recipients or transferability.

\section{Investment under heterogeneous beliefs}
\label{sec:investment}
\subsection{A minimal private decision framework}
An investor deciding whether to fund an agent faces uncertainty about its capabilities, the difficulty of a target, competitors, and contractual settlement. Let $\mathcal I_{i,t}$ be investor $i$'s information at time $t$, and let $Q_{i,t}$ be its subjective probability assessment. For a proposed action $a$, write
\begin{equation}
 \Pi_{i,t}(a)=\E_{Q_{i,t}}\!\left[
 R_i(a)-C_i^{\mathrm{net}}(a)\mid\mathcal I_{i,t}\right].
 \label{eq:subjective}
\end{equation}
Here $R_i$ is the investor's contracted share of payments, and $C_i^{\mathrm{net}}$ includes its resource purchases, fees, audit costs, and financing costs, less support it is entitled to receive. The horizon and discounting convention must be specified in an application. The receipt of a subsidy should be accounted for once, either in net cost or in receipts. Equation~\eqref{eq:subjective} describes a private assessment, not a correct population forecast.

There is no common-prior, calibration, or rational-expectations assumption. Investors can have different information and different models even when they observe the same public record. The emphasis on heterogeneous assessments is consistent with \citet{harrison}; their specific results about resale speculation also depend on their market restrictions and are not assumed here. An investor may value a claim because it expects a payout, because it expects to sell the claim, or for both reasons. A nontransferable financing contract removes the resale channel while leaving disagreement about research outcomes intact.

A risk-neutral participant with sufficient liquidity might select the feasible action with the highest subjective net value, including the option of doing nothing. Risk aversion, exposure limits, grant obligations, and scientific preferences can change that rule. An agent acting for an investor should inherit these constraints. The feasible set can include continuing a run, trying another strategy, funding a different team, purchasing information, or stopping. Repeated decisions allow models to update, but updating does not ensure eventual agreement or accuracy.

A simple illustration makes the distinction concrete. Suppose a claim pays 100 units if its team obtains an eligible award, and an attempt costs 20 net units. An investor assigning probability 0.30 to that event sees expected net receipts of 10; another assigning 0.10 sees a loss of 10. These figures are assessments of the same opportunity, not conflicting measurements of an observable constant. They also ignore risk and alternative uses of capital. If a sharing rule makes the conditional payment uncertain, each investor needs a distribution over its eventual share as well as a probability of producing a valid proof.

\subsection{Different questions require different forecasts}
At least four beliefs matter. A participant may assess whether a conjecture is true, whether a particular procedure can produce an admissible proof within a budget, whether its submission will qualify for an award, and whether the resulting receipts will exceed its costs. These events are not interchangeable. A true conjecture may be inaccessible to current tools. A successful formal proof may fail a timing or attribution rule. A paid result may still be privately unprofitable. Conversely, a valuable refutation may be rewarded under a contract that explicitly admits it.

Forecasts can use expert judgments, local experiments, historical traces, current platform activity, and market prices. The investor chooses how much weight to assign to each signal and whether additional analysis is worth its cost. An expert who understands an area may identify a hidden obstruction absent from a language model's estimate. A local agent may have better evidence about its own prover configuration. Disagreement is informative about the state of knowledge and should not be erased by publishing a single number without its event definition, horizon, and assumptions.

A prediction contract adds another layer. Its traded price concerns a specified settlement event and reflects the behavior of market participants under the contract's rules. It is not automatically a probability for an investor's own proof search. Liquidity, incentives to acquire information, trading constraints, and the ability of traders to affect outcomes matter. \citet{grossman} explains why costly information and perfectly revealing prices are difficult to reconcile in their model. The implication for this framework is modest: public signals need not eliminate the value of independent research, and their informativeness must be evaluated rather than presumed.

\subsection{Private return and the allocation of scientific effort}
Private profitability need not rank tasks by scientific value. An easy, heavily funded lemma can dominate an important but underfunded question in an investor's objective. A reusable tool can have wide benefits that no single mission pays for. Sponsors therefore influence allocation not only by the size of terminal prizes but also by funding infrastructure, broad exploration, and intermediate contributions. The sponsor's own assessment of value may also be uncertain, purpose-dependent, and contested.

Competition makes private returns endogenous. A public signal that attracts many agents can reduce each entrant's chance of obtaining an exclusive award or dilute a shared pool. The same signal can improve collective completion probability while reducing returns to late entrants. Search costs are correlated when participants use similar models and libraries, and successes on related lemmas can jointly change several tasks' prospects. Diversifying across task labels may offer little real diversification. These interactions are reasons to study entry and portfolio choice, not grounds for presenting a single return estimate as the market's objective opportunity set.

A funding system should record enough information to distinguish spending induced by a reward from spending that would have occurred anyway. At the same time, disclosure rules must account for strategic delay. A team may withhold a useful partial result if early publication benefits competitors more than its own claim. Intermediate awards, precommitted team shares, and narrowly limited priority protection can reduce that tension. None universally eliminates it; their effect depends on the contribution structure and on what can be credibly observed.

\section{Coordination across platforms}
\label{sec:coordination}
\subsection{Shared task identity and the meaning of time}
Multiple coordinators and verifiers can broaden access, create redundancy, and let participants choose services suited to their work. They also create a coordination problem. Two platforms may display the same mission under different identifiers, run different library versions, or disagree about when a submission became eligible. An interoperable market needs common task identifiers, versioned specifications, authenticated submissions, and a rule for reconciling records. Participation through a different interface should not silently change the underlying contract.

Three times should be distinguished: the time a platform receives a submission, the time a commitment becomes part of the agreed record, and the time verification finishes. None is necessarily the time of discovery. A slow verifier should not erase an otherwise timely commitment if the contract permits later verification. Conversely, a participant's self-reported local clock should not establish priority. A common record can order events without establishing exact physical time, and even timestamp authorities rely on explicit trust and clock assumptions.

A blockchain is one possible shared record. Its value is strongest when parties need a common account of commitments and settlement without delegating all control to one platform. Finality rules, inclusion delays, validator incentives, and outages still affect the observed order. Signed append-only logs, independent witnesses, or a trusted timestamp authority are alternatives; RFC~3161 specifies a timestamp service for evidence that data existed before a particular time \citep{rfc3161}. None of these mechanisms proves who independently discovered the result. The reward contract should state which record is authoritative and how uncertainty, outages, and migration are handled.

\subsection{Priority, nearby submissions, and fixed reward pools}
Exclusive payment to the first valid submission is a useful baseline, but it is not logically implied by formal verification. A sponsor can instead reward a set of nearby valid commitments, or divide funds among distinct contributions. Prize concentration strengthens the gain from finishing ahead of others; sharing can reduce exposure to small differences in latency and recognize parallel work. \citet{moldovanu} shows that single versus multiple prizes can be optimal under different cost conditions in a contest maximizing aggregate effort. Their result does not establish an optimal time window for proof production or an optimal measure of scientific progress.

For concreteness, consider a round with a fixed commitment deadline $T_c$ and a later common reveal deadline $T_r$. Commitments bind the participant account, task version, artifact, and a nonce. No eligible artifact is publicly revealed through this procedure until the commitment phase closes. Verification can continue after revelation under a published processing and appeal schedule. Only commitments opened correctly by $T_r$ and eventually accepted under the agreed procedure enter the reward calculation. An earlier invalid or unopened commitment therefore cannot move the reference time.

Let $\tau_a$ be the canonical record time of commitment $a$, let $\mathcal V$ be the nonempty set of valid, timely revealed, eligible commitments, and let $\tau_* =\min_{a\in\mathcal V}\tau_a$. For an announced window $\Delta\geq0$, define
\[
 \mathcal W_\Delta=\{a\in\mathcal V:\tau_a\leq\tau_*+\Delta\},
 \qquad P_a=\frac{B^T}{|\mathcal W_\Delta|}\I\{a\in\mathcal W_\Delta\}.
\]
If there is no valid commitment, the pool is retained or returned according to the contract. This equal-sharing rule is an illustration, not an optimal mechanism. The round schedule, window, budget, and eligibility unit must be fixed before participation. A strict first-winner rule requires an explicit tie-breaker when recorded times tie; a zero window alone does not resolve ties.

The shared reveal phase limits copying from disclosures made through the protocol while commitments remain open. Binding an account to the commitment prevents someone from claiming another account's unopened commitment as its own. These protections do not establish independent authorship: participants can coordinate privately, copy an already public proof, or create several accounts. A contract needs a declared unit of entitlement, duplicate-claim rules, and an identity or affiliation policy appropriate to its reward size. Small bonds or admission limits can deter cheap spam, but also obstruct entry. Deduplicating identical files is insufficient because superficial changes can produce many different proof artifacts.

A wide window may recognize near-simultaneous work but dilute incentives and reward redundant effort. A narrow window leaves participants exposed to network delay. Rounds facilitate a clean commit--reveal sequence but delay disclosure and settlement. Continuous operation can use successive rounds, yet must specify how already published proofs affect later eligibility. Where ordering is unreliable, a preannounced coarser batch rule is more defensible than a claim of spurious timestamp precision. Calibration of these choices is a research problem requiring latency and participation data.

\subsection{Verification, settlement, and governance}
Platforms should exchange evidence bundles rather than unverifiable declarations of success. A bundle can contain the canonical task version, commitment receipt, disclosed artifact, exact checking environment, verification report, and payout identifier. Multiple services can independently check the same artifact. Counting their votes is not automatically meaningful when they share code or infrastructure; the benefits of redundancy depend on independent failure modes. Conflicting reports should trigger a specified reconciliation procedure and re-execution against the authoritative environment.

Settlement must also prevent duplicate withdrawal from a common pool. The same proof may legitimately receive distinct awards from distinct sponsors, but one entitlement must not be redeemed twice merely because two front ends accepted it. Each award needs a globally reconciled redemption state or an agreed settlement authority. A ledger can enforce this state transition, while off-chain services supply proof and audit evidence. Escrow sufficiency, verifier authorization, and dispute timing must be established before automatic payment is credible.

Governance determines who may change a task, approve a new dependency, suspend an unsafe verifier, or resolve an attribution dispute. It also determines how conflicts are handled when a platform both invests in a solver and certifies its result. Public rules, disclosed affiliations, independent review for contested decisions, and separation of responsibilities can constrain such conflicts. Decentralization is useful only insofar as the resulting arrangement makes these powers more contestable or reliable. Replacing a platform operator with an opaque committee or concentrated token vote does not by itself accomplish that goal.

\section{Institutional implementations and emerging examples}
\label{sec:implementations}
\subsection{Relating institutional practice to the framework}
The roles above provide a way to compare institutions as their functions and scope evolve. The examples below identify practices that already address parts of the economic problem and the additional arrangements needed to connect them. Four questions organize the comparison: what objective or service is offered, which evidence is accepted, how resources and rewards are committed, and who has authority over disputed decisions. These questions connect institutional experience to the assumptions of the abstract framework.

Descriptions of existing capabilities are based on the cited public primary materials available on October 1, 2026. Proposed extensions are identified separately. The examples do not establish a single integrated market. Evidence about operating costs, participation, verification, and disputed outcomes would help determine which distinctions in the framework matter in practice and where the framework needs refinement.

\subsection{Demand, sponsorship, and programmable payment}
Government programs and corporate verification efforts illustrate the demand side. DARPA's expMath program seeks AI capabilities that assist mathematics, including problem decomposition and useful intermediate results \citep{darpa}. AWS describes applications of automated reasoning to the security properties of its services \citep{aws}. These examples establish motivations for mathematical development and formal assurance. They do not show that either institution operates an open, token-funded proof market. Procurement and grant programs would need their own task, disclosure, and contracting arrangements before connecting to one.

The Justin Sun Prize provides an example closer to an outcome reward. Its published rules describe cryptocurrency prizes, formal verification, and recognition of solving and formalizing contributions \citep{sun}. They also retain an adjudication process. A machine-checked artifact is consequently part of the award evidence, rather than a sufficient description of an automatic payment rule. The case illustrates how a sponsor can value a mathematical milestone without simultaneously supplying a distributed work platform or underwriting every participant's unsuccessful attempts.

Programmable payment protocols can make the interfaces between these roles easier to operate. For example, x402 describes a protocol for requesting and making payments through HTTP interactions \citep{x402}. Such an interface could let an authorized agent purchase services or receive a small payment under a spending policy. Faster execution is a possible benefit when it reduces repeated administrative intervention; it depends on the payment rail, finality, fees, and integration. The protocol does not certify compute delivery, create a theorem reward, or require this ecosystem to issue its own currency. Restricted wallets and explicit authority to spend remain distinct from the proof-search algorithm.

\subsection{OpenMath: tasks, commitments, and rewards}
OpenMath's official agent tooling describes discovery of open theorem tasks, local Lean or Rocq proof work, commitment and revelation of submissions, and reward claims \citep{openmath}. It therefore illustrates a bundle of task discovery, submission recording, and settlement interfaces. A commitment mechanism is particularly relevant to the separation between registering an artifact and disclosing it. Earlier work by \citet{lim} also studies blockchain incentives for collaborative formalization, so the broad idea of combining formal proofs with a ledger is not new.

This workflow provides a concrete setting for examining the coordination questions in Section~\ref{sec:coordination}. Recognition across platforms would require agreement on task versions, eligibility, and the treatment of conflicting records. Exploratory support would additionally require expenditure evidence and a reserved budget; sharing among nearby submissions would require an announced entitlement and timing rule. The cited tooling describes the submission and reward workflow, while these additional arrangements remain design possibilities in the present analysis. Experience with actual commitments, verification delays, and reward claims could inform their design.

\subsection{Prove2Me: missions and compositional progress}
\citet{prove2me} describes a collaborative platform in which human participants bring agents to mathematical formalization missions. Its architecture distinguishes statements from proofs, uses human oversight of canonical mission content, and supports proof-sketches whose dependencies can be pursued separately. A root statement can thus have a checked reduction while remaining conditional on unresolved supporting statements. Reusable formal results and contribution credit are part of the coordination design. These capabilities correspond principally to curation, distributed execution, and the organization of verification evidence.

This architecture is a natural setting in which to examine the intermediate rewards in Section~\ref{sec:funding}. A sponsor could reserve awards for reviewed dependencies and support participants' compute use while those tasks are open. Such an extension would need a sponsor budget, expenditure evidence, eligibility rules, and a settlement interface. Contribution credit in the cited platform describes recognition of work; a transferable payment or enforceable reward claim carries additional financial obligations. Experience with mission decomposition, dependency closure, and actual reuse could help determine which contributions are suitable for funded contracts and how their value should be assessed.

\subsection{Forecasts, expert information, and investor decisions}
Kalshi's filed Millennium Prize contract terms identify an award announcement by the Clay Mathematics Institute as the relevant settlement event \citep{kalshi}. That event differs from the truth of a conjecture, the production of a Lean proof by a particular team, or that team's receipt of another platform's bounty. A price for the filed event could inform an investor's broader assessment of progress, but translating it into a private success probability requires additional judgment. The example illustrates the importance of reading the actual contract. It does not supply a current price, establish liquidity, or show that the event contract is calibrated to an agent's research horizon.

Information providers may supply expert assessments, forecasts, or evidence from local proof experiments. In the framework of Section~\ref{sec:investment}, investors decide how these inputs enter their own estimates of returns and costs under a stated task, budget, and horizon. Evaluating this interface requires recording forecasts before outcomes are known and including failed and discontinued attempts, as discussed in Section~\ref{sec:research}. These requirements connect the information-provider role to the allocation decisions and comparative experiments developed in this paper.

Research prototypes provide complementary evidence about automated allocation. \citet{hunt} studies a simulated bounty environment for collaborative autoformalization, including agent-generated tasks and checked results. \citet{costquality} studies decisions that change the cost and success of agentic Lean proving. These works motivate experiments with task selection and continuation policies. They do not establish a profitable autonomous investment fund or the performance of a live cross-platform reward economy. Lean's validation documentation further emphasizes the care required to establish exactly which statement, assumptions, and checking path a purported proof validates \citep{lean}.

\section{Risks and research agenda}
\label{sec:research}
\subsection{Contractible output does not eliminate hidden action}
Formal acceptance makes one dimension of output unusually tractable. It leaves hidden actions in research effort, model choice, strategic disclosure, provider billing, and account creation. Paying for metered activity can induce waste; paying only for completion can suppress exploration. Intermediate awards can finance a useful library or stimulate artificial fragmentation. The appropriate comparison is therefore between feasible institutional arrangements, including their monitoring costs and exclusion effects. A more elaborate protocol is not an improvement unless the additional evidence changes behavior in a useful way.

Statement quality is an especially consequential risk. A sponsor can pay promptly for a formally valid result that answers the wrong question, relies on unintended assumptions, or proves a weaker claim than participants expect. The remedy begins before search: independent review of the canonical statement, dependency policy, and acceptance environment. For high-value tasks, a pilot submission and adversarial validation exercise can reveal weaknesses before large sums are committed. Later corrections require versioning and compensation rules; silently replacing a target after work has begun undermines commitment.

\subsection{Concentration, manipulation, and correlated mistakes}
Compute subsidies can broaden entry while concentrating activity in the providers whose receipts the platform trusts. Large investors may benefit from cheaper compute, private information, and diversified portfolios. Anti-Sybil requirements can in turn favor well-established participants. These are institutional tradeoffs, not arguments that openness or identity controls are inherently sufficient. Evaluation should examine who enters, who receives funding, and how awards are distributed, alongside completion rates.

Forecasts and rewards can interact strategically. An actor who trades on a milestone and also performs the work can influence the settlement event. An information provider can benefit from directing investors toward affiliated teams, and a verifier can have a financial stake in a contested award. Research on prediction markets with outcome-influencing participants makes clear why the information and action channels cannot always be analyzed separately \citep{manipulation}. Disclosure, position limits where appropriate, independent adjudication, and alternative elicitation designs are candidates for study. No single measure guarantees truthful forecasts.

Shared models can produce correlated errors in both search and valuation. If many agents overestimate the same route, a market can direct a large amount of computation into the same failure. Token appreciation can further detach attention from the value of the underlying proof services. Monetary accounts should distinguish sponsor-funded payouts, secondary transfers, and newly issued claims. More tokens do not create more model capacity, and an internally balanced ledger does not demonstrate backing, purchasing power, or solvency.

\subsection{A program of comparative evaluation}
The first research priority is a controlled comparison of funding mechanisms. Missions can be assigned comparable total budgets with different allocations to advance credits, capped reimbursement, intermediate awards, and terminal prizes. Outcomes should include accepted and reused artifacts, time to completion, total compute and human review cost, audit failures, and participation by resource-constrained teams. Disclosure and benchmark contamination must be controlled. Because work on one mission can benefit another, spillovers and shared dependencies should be recorded rather than treated as independent observations.

A second priority is the interaction of entry, sharing, and latency. Experiments can vary reward windows and disclosure schedules while measuring duplicate effort, strategic delay, identity abuse, and the rate of disputed settlements. The central question is whether sharing preserves enough incentive to finish while reducing arbitrary advantages from infrastructure. A third priority is forecast evaluation: cost and success estimates should be recorded before attempts, include unsuccessful and censored runs, and be compared against simple baselines. Calibration of one event does not establish calibration of a different payoff claim.

Dynamic models can subsequently examine learning, congestion, portfolio choice, and competition among platforms. Mean-field or network models may become useful when their state variables and interaction assumptions are supported by data. They are not needed to define the roles or justify an accounting identity, and no aggregate equilibrium is asserted here. The immediate contribution of the framework is to identify the evidence, contractual choices, and outcomes such models would need to explain.

\section{Conclusion}
AI-assisted search and formal proof checking create the conditions for a different way to organize mathematical work. Human effort and computation can be committed to uncertain searches, and successful formal results can then be checked and reused by others. The economic problem is how to connect demand for those results to the resources consumed along the way. Formal acceptance supplies important evidence, while financing, contribution attribution, and payment require their own rules.

A functional account makes those rules comparable across institutions. It distinguishes sponsors who value knowledge from investors who assess contingent returns, and connects their decisions to coordination, proof work, verification, information, and settlement. Compute support, intermediate awards, and sharing among nearby submissions address different parts of the funding problem. Each has incentive and governance costs that depend on the objective, the evidence available, and the distribution of authority. Heterogeneous beliefs and competition further shape how participants respond.

Existing initiatives supply concrete practices against which this framework can be examined and revised. Comparisons of their workflows and outcomes can reveal which contractual distinctions are useful, which assumptions are too restrictive, and which interfaces would enable further cooperation. Bounded experiments with explicit budgets and acceptance rules can then test whether particular arrangements improve participation and produce more checked, reusable mathematical progress per unit of resources. That relationship between institutional experience and mechanism analysis is the basis for developing the framework further.

\clearpage
\begin{thebibliography}{99}
\small
\bibitem[Adams et~al.(2001)]{rfc3161} Adams, Carlisle, Pat Cain, Denis Pinkas, and Robert Zuccherato. 2001. ``Internet X.509 Public Key Infrastructure Time-Stamp Protocol (TSP).'' \href{https://www.rfc-editor.org/rfc/rfc3161}{RFC 3161}.
\bibitem[Aghion and Tirole(1994)]{aghiontirole} Aghion, Philippe, and Jean Tirole. 1994. ``The Management of Innovation.'' \emph{Quarterly Journal of Economics} 109(4): 1185--1209. \href{https://doi.org/10.2307/2118360}{doi:10.2307/2118360}.
\bibitem[AWS(n.d.)]{aws} Amazon Web Services. n.d. ``Provable Security.'' \href{https://aws.amazon.com/security/provable-security/}{Official description}. Accessed October 1, 2026.
\bibitem[Brown et~al.(2026)]{hunt} Brown, Chad E., Cezary Kaliszyk, and Josef Urban. 2026. ``Agent Hunt: Bounty Based Collaborative Autoformalization With LLM Agents.'' \href{https://arxiv.org/abs/2603.06737}{arXiv:2603.06737}.
\bibitem[Chakraborty and Das(2014)]{manipulation} Chakraborty, Mithun, and Sanmay Das. 2014. ``On Manipulation in Prediction Markets When Participants Influence Outcomes Directly.'' \href{https://arxiv.org/abs/1407.7015}{arXiv:1407.7015}.
\bibitem[Chen et~al.(2026)]{prove2me} Chen, Shuze, Kunal Marwaha, Xiaoyang Lu, Henry Yuen, and Tianyi Peng. 2026. ``Prove2Me: An Open Collaborative Platform for Scaling Math Formalization.'' \href{https://arxiv.org/html/2608.28433v2}{arXiv:2608.28433v2}.
\bibitem[DARPA(n.d.)]{darpa} DARPA. n.d. ``expMath: Exponential Mathematics.'' \href{https://www.darpa.mil/research/programs/expmath-exponential-mathematics}{Official program description}. Accessed October 1, 2026.
\bibitem[Grossman and Stiglitz(1980)]{grossman} Grossman, Sanford J., and Joseph E. Stiglitz. 1980. ``On the Impossibility of Informationally Efficient Markets.'' \emph{American Economic Review} 70(3): 393--408. \href{https://www.jstor.org/stable/1805228}{JSTOR stable record}.
\bibitem[Harrison and Kreps(1978)]{harrison} Harrison, J. Michael, and David M. Kreps. 1978. ``Speculative Investor Behavior in a Stock Market with Heterogeneous Expectations.'' \emph{Quarterly Journal of Economics} 92(2): 323--336. \href{https://doi.org/10.2307/1884166}{doi:10.2307/1884166}.
\bibitem[KalshiEX(2025)]{kalshi} KalshiEX LLC. 2025. ``Will a Millennium Prize Be Awarded Before $\langle$date$\rangle$?'' Contract terms in \href{https://www.cftc.gov/filings/ptc/ptc08142528943.pdf}{CFTC filing}, August 14, pp.~7--8.
\bibitem[Kremer(1998)]{kremer} Kremer, Michael. 1998. ``Patent Buyouts: A Mechanism for Encouraging Innovation.'' \emph{Quarterly Journal of Economics} 113(4): 1137--1167. \href{https://doi.org/10.1162/003355398555865}{doi:10.1162/003355398555865}. Locally archived open version: \href{https://www.nber.org/system/files/working_papers/w6304/w6304.pdf}{NBER Working Paper 6304}, December 1997.
\bibitem[Lean Project(n.d.)]{lean} Lean Project. n.d. ``Validating a Lean Proof.'' \href{https://lean-lang.org/doc/reference/latest/ValidatingProofs/}{Lean Language Reference}. Accessed October 1, 2026.
\bibitem[Lim et~al.(2021)]{lim} Lim, Jin Xing, Barnab\'e Monnot, Shaowei Lin, and Georgios Piliouras. 2021. ``A Blockchain-Based Approach for Collaborative Formalization of Mathematics and Programs.'' \href{https://arxiv.org/abs/2111.10824}{arXiv:2111.10824}.
\bibitem[Manso(2011)]{manso} Manso, Gustavo. 2011. ``Motivating Innovation.'' \emph{Journal of Finance} 66(5): 1823--1860. \href{https://doi.org/10.1111/j.1540-6261.2011.01688.x}{doi:10.1111/j.1540-6261.2011.01688.x}. Locally archived author manuscript dated November 26, 2010.
\bibitem[Moldovanu and Sela(2001)]{moldovanu} Moldovanu, Benny, and Aner Sela. 2001. ``The Optimal Allocation of Prizes in Contests.'' \emph{American Economic Review} 91(3): 542--558. \href{https://doi.org/10.1257/aer.91.3.542}{doi:10.1257/aer.91.3.542}. Locally archived author manuscript dated June 5, 2000.
\bibitem[R\"ognvaldsson et~al.(2026)]{costquality} R\"ognvaldsson, K\'ari, Chenhao Sun, Jasper Dekoninck, and Martin Vechev. 2026. ``Optimizing the Cost-Quality Tradeoff of Agentic Theorem Provers in Lean.'' \href{https://arxiv.org/abs/2606.04883}{arXiv:2606.04883}.
\bibitem[Shentu Foundation(n.d.)]{openmath} Shentu Foundation. n.d. \emph{OpenMath Skills}. \href{https://github.com/shentufoundation/openmath-skills}{Official repository}, including theorem discovery and submission workflows. Accessed October 1, 2026.
\bibitem[The Justin Sun Prize(2026)]{sun} The Justin Sun Prize. 2026. \emph{Selection Rules}, version 1.0, effective September 16. \href{https://www.hejustinsun.com/prize/rules}{Official rules}. Accessed October 1, 2026.
\bibitem[x402 Foundation(n.d.)]{x402} x402 Foundation. n.d. \emph{x402: A Payments Protocol for the Internet}. \href{https://github.com/x402-foundation/x402}{Official repository}. Accessed October 1, 2026.
\end{thebibliography}
\end{document}
